%------------------------------------------------------------------------------------
%  TABLEAU PAYSAGE LARGE (pdflscape + booktabs + tabularx)
%------------------------------------------------------------------------------------
\begin{landscape}
\begin{table}[p]
  \centering
  \small
  \caption{Matrice d'évaluation comparative multicritères des architectures logicielles et typographiques}
  \label{tab:evaluation_paysage}
  \begin{tabularx}{\linewidth}{p{3.5cm} p{3.5cm} c c c c X}
    \toprule
    \textbf{Moteur / Chaîne} & \textbf{Standard de Polices} & \textbf{UTF-8 Natif} & \textbf{Microtypo} & \textbf{Vitesse} & \textbf{Complexité} & \textbf{Conclusion Opérationnelle} \\
    \midrule
    pdf\LaTeX{} classique & T1 / Type 1 (8 bits) & Non (inputenc) & Limitée & Rapide & Élevée (bidouilles) & Obsolète pour nouveaux projets académiques. \\
    \addlinespace
    Xe\LaTeX{} & OpenType / TrueType & Oui & Basique & Moyenne & Moyenne & Préféré historiquement pour langues non latines. \\
    \addlinespace
    \textbf{Lua\LaTeX{} moderne} & \textbf{OpenType / TrueType} & \textbf{Oui (100\,\%)} & \textbf{Complète} & \textbf{Excellente} & \textbf{Épurée} & \textbf{Standard contemporain recommandé par l'UCLouvain.} \\
    \midrule
    \multicolumn{7}{l}{\textbf{Bilan comparatif des fonctionnalités de publication :}} \\
    \addlinespace
    Classe KOMA \texttt{scrbook} & Système KOMA-Script & Oui & Oui & -- & Modérée & Socle de base solide mais sans encadrements prêts à l'emploi. \\
    \addlinespace
    \texttt{ClasseGuillaume} & \textbf{Architecture Modulaire} & \textbf{Oui} & \textbf{Pénalités max} & -- & \textbf{Immédiate} & \textbf{Intègre \texttt{\textbackslash todoH}, 3 géométries, logo scrlayer et booktabs.} \\
    \bottomrule
  \end{tabularx}
  \vspace{4mm}
  \begin{minipage}{\linewidth}
    \footnotesize
    \textit{Note méthodologique :} Ce tableau panoramique démontre l'utilisation conjointe de l'environnement \texttt{landscape} fourni par le paquet \textsf{pdflscape} et de la classe de document. La rotation est physique dans le fichier PDF généré, garantissant une consultation parfaite sans devoir tourner manuellement l'écran.
  \end{minipage}
\end{table}
\end{landscape}

